% Mathematics and Physics statements
%
% The following boxes are provided:
%   Mathematical statements:
%   Axiom:          \axiomr{name}{axiom:reference label}{statement}
%   Definition:     \defnr{name}{defn:reference label}{statement}
%   Theorem:        \thmr{name}{thm:reference label}{statement}}
%   Proposition:    \proppfr{name}{prop:reference label}{statement}{proof}
%   Lemma:          \lempfr{name}{lem:reference label}{statement}{proof}
%   Corollary:      \corpfr{name}{cor:reference label}{statement}{proof}
%   Conjecture:     \conjr{name}{conj:reference label}{statement}

%   Physics statements:
%   Postulate:      \postr{name}{post:reference label}{statement}
%   Terminology:    \tmgr{name}{tmg:reference label}{statement}
%   Principle:      \prr{name}{pr:reference label}{statement}
%   Fact:           \factpfr{name}{fact:reference label}{statement}{proof}
%   Result:         \respfr{name}{res:reference label}{statement}{proof}


%   Notations
%   Note:           \nt{statement}

%-------------------------------Mathematical statements---------------------------------

%Axiom
\newtcbtheorem[number within=section]{myaxiom}{Axiom}
{
    enhanced,
    frame hidden,
    titlerule=0mm,
    toptitle=1mm,
    bottomtitle=1mm,
    fonttitle=\bfseries\large,
    coltitle=black,
    colbacktitle=gray!20!white,
    colback=gray!10!white,
}{axiom}

\NewDocumentCommand{\axiom}{m+m}{
    \begin{myaxiom}{#1}{}
        #2
    \end{myaxiom}
}

\NewDocumentCommand{\axiomr}{mm+m}{
    \begin{myaxiom}{#1}{#2}
        #3
    \end{myaxiom}
}

% Definition
\newtcbtheorem[use counter from=myaxiom]{mydefinition}{Definition}
{
    enhanced,
    frame hidden,
    titlerule=0mm,
    toptitle=1mm,
    bottomtitle=1mm,
    fonttitle=\bfseries\large,
    coltitle=black,
    colbacktitle=red!20!white,
    colback=red!10!white,
}{defn}

\NewDocumentCommand{\defn}{m+m}{
    \begin{mydefinition}{#1}{}
        #2
    \end{mydefinition}
}

\NewDocumentCommand{\defnr}{mm+m}{
    \begin{mydefinition}{#1}{#2}
        #3
    \end{mydefinition}
}

% Theorem
\newtcbtheorem[use counter from=myaxiom]{mytheorem}{Theorem}
{
    enhanced,
    frame hidden,
    titlerule=0mm,
    toptitle=1mm,
    bottomtitle=1mm,
    fonttitle=\bfseries\large,
    coltitle=black,
    colbacktitle=cyan!20!white,
    colback=cyan!10!white,
}{thm}

\NewDocumentCommand{\thm}{m+m}{
    \begin{mytheorem}{#1}{}
        #2
    \end{mytheorem}
}

\NewDocumentCommand{\thmr}{mm+m}{
    \begin{mytheorem}{#1}{#2}
        #3
    \end{mytheorem}
}

\newenvironment{theoremproof}{
    {\noindent{\it \textbf{proof of the theorem.}}}
    \tcolorbox[blanker, breakable, left=5mm, right=1mm, parbox=false,
        before upper={\parindent15pt},
        after skip=10pt,
        borderline west={1mm}{0pt}{cyan!20!white}]
}{
    \textcolor{cyan!20!white}{\hbox{}\nobreak\hfill$\blacksquare$} 
    \endtcolorbox
}

\NewDocumentCommand{\thmpf}{m+m+m}{
    \begin{mytheorem}{#1}{}
        #2
    \end{mytheorem}

    \begin{theoremproof}
        #3
    \end{theoremproof}
}

\NewDocumentCommand{\thmpfr}{mm+m+m}{
    \begin{mytheorem}{#1}{#2}
        #3
    \end{mytheorem}

    \begin{theoremproof}
        #4
    \end{theoremproof}
}

% Proposition
\newtcbtheorem[use counter from=myaxiom]{myproposition}{Proposition}
{
    enhanced,
    frame hidden,
    titlerule=0mm,
    toptitle=1mm,
    bottomtitle=1mm,
    fonttitle=\bfseries\large,
    coltitle=black,
    colbacktitle=blue!15!white,
    colback=blue!10!white,
}{prop}

\NewDocumentCommand{\prop}{m+m}{
    \begin{myproposition}{#1}{}
        #2
    \end{myproposition}
}

\NewDocumentCommand{\propr}{mm+m}{
    \begin{myproposition}{#1}{#2}
        #3
    \end{myproposition}
}

\newenvironment{propositionproof}{
    {\noindent{\it \textbf{proof of the proposition.}}}
    \tcolorbox[blanker, breakable, left=5mm, right=1mm, parbox=false,
        before upper={\parindent15pt},
        after skip=10pt,
        borderline west={1mm}{0pt}{blue!15!white}]
}{
    \textcolor{blue!15!white}{\hbox{}\nobreak\hfill$\blacksquare$} 
    \endtcolorbox
}

\NewDocumentCommand{\proppf}{m+m+m}{
    \begin{myproposition}{#1}{}
        #2
    \end{myproposition}

    \begin{propositionproof}
        #3
    \end{propositionproof}
}

\NewDocumentCommand{\proppfr}{mm+m+m}{
    \begin{myproposition}{#1}{#2}
        #3
    \end{myproposition}

    \begin{propositionproof}
        #4
    \end{propositionproof}
}

% Lemma
\newtcbtheorem[use counter from=myaxiom]{mylemma}{Lemma}
{
    enhanced,
    frame hidden,
    titlerule=0mm,
    toptitle=1mm,
    bottomtitle=1mm,
    fonttitle=\bfseries\large,
    coltitle=black,
    colbacktitle=green!20!white,
    colback=green!10!white,
}{lem}

\NewDocumentCommand{\lem}{m+m}{
    \begin{mylemma}{#1}{}
        #2
    \end{mylemma}
}

\NewDocumentCommand{\lemr}{mm+m}{
    \begin{mylemma}{#1}{#2}
        #3
    \end{mylemma}
}

\newenvironment{lemmaproof}{
    {\noindent{\it \textbf{proof of the lemma.}}}
    \tcolorbox[blanker, breakable, left=5mm, right=1mm, parbox=false,
        before upper={\parindent15pt},
        after skip=10pt,
        borderline west={1mm}{0pt}{green!20!white}]
}{
    \textcolor{green!20!white}{\hbox{}\nobreak\hfill$\blacksquare$} 
    \endtcolorbox
}

\NewDocumentCommand{\lempf}{m+m+m}{
    \begin{mylemma}{#1}{}
        #2
    \end{mylemma}

    \begin{lemmaproof}
        #3
    \end{lemmaproof}
}

\NewDocumentCommand{\lempfr}{mm+m+m}{
    \begin{mylemma}{#1}{#2}
        #3
    \end{mylemma}

    \begin{lemmaproof}
        #4
    \end{lemmaproof}
}

% Corollary
\newtcbtheorem[use counter from=myaxiom]{mycorollary}{Corollary}
{
    enhanced,
    frame hidden,
    titlerule=0mm,
    toptitle=1mm,
    bottomtitle=1mm,
    fonttitle=\bfseries\large,
    coltitle=black,
    colbacktitle=cyan!15!white,
    colback=cyan!10!white,
}{cor}

\NewDocumentCommand{\cor}{m+m}{
    \begin{mycorollary}{#1}{}
        #2
    \end{mycorollary}
}

\NewDocumentCommand{\corr}{mm+m}{
    \begin{mycorollary}{#1}{#2}
        #3
    \end{mycorollary}
}

\newenvironment{corollaryproof}{
    {\noindent{\it \textbf{proof of the corollary.}}}
    \tcolorbox[blanker, breakable, left=5mm, right=1mm, parbox=false,
        before upper={\parindent15pt},
        after skip=10pt,
        borderline west={1mm}{0pt}{cyan!15!white}]
}{
    \textcolor{cyan!15!white}{\hbox{}\nobreak\hfill$\blacksquare$} 
    \endtcolorbox
}

\NewDocumentCommand{\corpf}{m+m+m}{
    \begin{mycorollary}{#1}{}
        #2
    \end{mycorollary}

    \begin{corollaryproof}
        #3
    \end{corollaryproof}
}

\NewDocumentCommand{\corpfr}{mm+m+m}{
    \begin{mycorollary}{#1}{#2}
        #3
    \end{mycorollary}

    \begin{corollaryproof}
        #4
    \end{corollaryproof}
}

% Conjecture
\newtcbtheorem[use counter from=myaxiom]{myconjecture}{Conjecture}
{
    enhanced,
    frame hidden,
    titlerule=0mm,
    toptitle=1mm,
    bottomtitle=1mm,
    fonttitle=\bfseries\large,
    coltitle=black,
    colbacktitle=yellow!20!white,
    colback=yellow!10!white,
}{conj}

\NewDocumentCommand{\conj}{m+m}{
    \begin{myconjecture}{#1}{}
        #2
    \end{myconjecture}
}

\NewDocumentCommand{\conjr}{mm+m}{
    \begin{myconjecture}{#1}{#2}
        #3
    \end{myconjecture}
}

%-------------------------------Physics statements---------------------------------
% Postulate
\newtcbtheorem[use counter from=myaxiom]{mypostulate}{Postulate}
{
    enhanced,
    frame hidden,
    titlerule=0mm,
    toptitle=1mm,
    bottomtitle=1mm,
    fonttitle=\bfseries\large,
    coltitle=black,
    colbacktitle=violet!20!white,
    colback=violet!10!white,
}{post}

\NewDocumentCommand{\post}{m+m}{
    \begin{mypostulate}{#1}{}
        #2
    \end{mypostulate}
}

\NewDocumentCommand{\postr}{mm+m}{
    \begin{mypostulate}{#1}{#2}
        #3
    \end{mypostulate}
}

% Terminology
\newtcbtheorem[use counter from=myaxiom]{myterminology}{Terminology}
{
    enhanced,
    frame hidden,
    titlerule=0mm,
    toptitle=1mm,
    bottomtitle=1mm,
    fonttitle=\bfseries\large,
    coltitle=black,
    colbacktitle=pink!50!white,
    colback=pink!25!white,
}{tmg}

\NewDocumentCommand{\tmg}{m+m}{
    \begin{myterminology}{#1}{}
        #2
    \end{myterminology}
}

\NewDocumentCommand{\tmgr}{mm+m+m}{
    \begin{myterminology}{#1}{#2}
        #3
    \end{myterminology}
}

% Principle
\newtcbtheorem[use counter from=myaxiom]{myprinciple}{Principle}
{
    enhanced,
    frame hidden,
    titlerule=0mm,
    toptitle=1mm,
    bottomtitle=1mm,
    fonttitle=\bfseries\large,
    coltitle=black,
    colbacktitle=teal!23!white,
    colback=teal!10!white,
}{pr}

\NewDocumentCommand{\pr}{m+m}{
    \begin{myprinciple}{#1}{}
        #2
    \end{myprinciple}
}

\NewDocumentCommand{\prr}{mm+m}{
    \begin{myprinciple}{#1}{#2}
        #3
    \end{myprinciple}
}

% Fact
\newtcbtheorem[use counter from=myaxiom]{myfact}{Fact}
{
    enhanced,
    frame hidden,
    titlerule=0mm,
    toptitle=1mm,
    bottomtitle=1mm,
    fonttitle=\bfseries\large,
    coltitle=black,
    colbacktitle=lime!25!white,
    colback=lime!11!white,
}{fact}

\NewDocumentCommand{\fact}{m+m}{
    \begin{myfact}{#1}{}
        #2
    \end{myfact}
}

\NewDocumentCommand{\factr}{mm+m}{
    \begin{myfact}{#1}{#2}
        #3
    \end{myfact}
}

\newenvironment{factproof}{
    {\noindent{\it \textbf{proof of the fact.}}}
    \tcolorbox[blanker, breakable, left=5mm, right=1mm, parbox=false,
        before upper={\parindent15pt},
        after skip=10pt,
        borderline west={1mm}{0pt}{lime!25!white}]
}{
    \textcolor{lime!25!white}{\hbox{}\nobreak\hfill$\blacksquare$} 
    \endtcolorbox
}

\NewDocumentCommand{\factpf}{m+m+m}{
    \begin{myfact}{#1}{}
        #2
    \end{myfact}

    \begin{factproof}
        #3
    \end{factproof}
}

\NewDocumentCommand{\factpfr}{mm+m+m}{
    \begin{myfact}{#1}{#2}
        #3
    \end{myfact}

    \begin{factproof}
        #4
    \end{factproof}
}

% Result
\newtcbtheorem[use counter from=myaxiom]{myresult}{Result}
{
    enhanced,
    frame hidden,
    titlerule=0mm,
    toptitle=1mm,
    bottomtitle=1mm,
    fonttitle=\bfseries\large,
    coltitle=black,
    colbacktitle=orange!20!white,
    colback=orange!10!white,
}{res}

\NewDocumentCommand{\res}{m+m}{
    \begin{myresult}{#1}{}
        #2
    \end{myresult}
}

\newenvironment{respf}{
    {\noindent{\it \textbf{Solution.}}}
    \tcolorbox[blanker, breakable, left=5mm, right=1mm, parbox=false,
        before upper={\parindent15pt},
        after skip=10pt,
        borderline west={1mm}{0pt}{orange!20!white}]
}{
    \textcolor{orange!20!white}{\hbox{}\nobreak\hfill$\blacksquare$} 
    \endtcolorbox
}

\NewDocumentCommand{\resp}{m+m+m}{
    \begin{myresult}{#1}{}
        #2
    \end{myresult}

    \begin{respf}
        #3
    \end{respf}
}

\NewDocumentCommand{\respfr}{mm+m+m}{
    \begin{myresult}{#1}{#2}
        #3
    \end{myresult}

    \begin{respf}
        #4
    \end{respf}
}

%-------------------------------Notations---------------------------------
% Note
\newenvironment{note}{
    {\noindent{\it \textbf{Note.}}}
    \tcolorbox[blanker, breakable, left=5mm, right=1mm, parbox=false,
        before upper={\parindent15pt},
        after skip=10pt,
        borderline west={1mm}{0pt}{blue!20!white}]
}{
    \textcolor{blue!20!white}{\hbox{}\nobreak\hfill$\blacksquare$} 
    \endtcolorbox
}

\NewDocumentCommand{\nt}{+m}{
    \begin{note}
        #1
    \end{note}
}
